\chapter{Implementazione}

\section{Stack Tecnologico}

La soluzione utilizza un moderno stack Python per data engineering:

\begin{table}[h]
\centering
\begin{tabular}{@{}llp{6cm}@{}}
\toprule
\textbf{Libreria} & \textbf{Versione} & \textbf{Utilizzo} \\
\midrule
\texttt{lxml} & 5.x & Parsing XML streaming \\
\texttt{polars} & 1.x & DataFrame processing \\
\texttt{pyarrow} & 18.x & Schema definition e Parquet I/O \\
\texttt{duckdb} & 1.x & Query engine SQL \\
\texttt{click} & 8.x & CLI framework \\
\texttt{tqdm} & 4.x & Progress bar \\
\bottomrule
\end{tabular}
\caption{Dipendenze principali}
\end{table}

\section{Sanitizzazione XML}

Un problema critico dei dataset RNA è la presenza di \textbf{caratteri XML non validi}. Questi includono:

\begin{itemize}
    \item Caratteri di controllo ASCII: \texttt{0x00-0x08}, \texttt{0x0B-0x0C}, \texttt{0x0E-0x1F}
    \item Entità numeriche malformate: \texttt{\&\#0;} - \texttt{\&\#31;}
\end{itemize}

La classe \texttt{CleanFileInputStream} risolve questo problema tramite un wrapper file-like che filtra lo stream in tempo reale:

\begin{lstlisting}[language=Python, caption=Sanitizzazione streaming]
class CleanFileInputStream:
    def __init__(self, filename):
        self.f = open(filename, 'rb')
        self.invalid_xml_chars = re.compile(
            b'[\x00-\x08\x0b\x0c\x0e-\x1f]|'
            b'&#(?:[0-8]|1[1-2]|1[4-9]|2[0-9]|3[0-1]);'
        )

    def read(self, size=-1):
        data = self.f.read(size)
        return self.invalid_xml_chars.sub(b'', data)
\end{lstlisting}

\begin{tipbox}[Design Choice]
Il pattern wrapper consente di integrare la sanitizzazione senza modificare il codice di parsing. Il parser \texttt{lxml} riceve uno stream già pulito.
\end{tipbox}

\section{Parsing Streaming}

Per gestire file multi-gigabyte senza esaurire la memoria, utilizziamo \texttt{lxml.etree.iterparse()} con cleanup aggressivo:

\begin{lstlisting}[language=Python, caption=Iterparse con memory management]
context = etree.iterparse(
    clean_stream, 
    events=("end",), 
    tag=f"{NS}AIUTO"
)

for event, elem in context:
    # Estrazione dati...
    
    # Release memory
    elem.clear()
    while elem.getprevious() is not None:
        del elem.getparent()[0]
\end{lstlisting}

Il pattern \texttt{elem.clear()} + deletion dei siblings previene memory leak tipici di \texttt{iterparse}.

\section{Parallelismo Multi-Processo}

L'elaborazione parallela sfrutta \texttt{ProcessPoolExecutor} per bypassare il GIL:

\begin{lstlisting}[language=Python, caption=Orchestrazione parallela]
with ProcessPoolExecutor(max_workers=workers) as executor:
    futures = {
        executor.submit(process_file, f, output): f 
        for f in files
    }
    
    for future in as_completed(futures):
        stats = future.result()
        # Aggregate statistics...
\end{lstlisting}

\begin{figure}[h]
\centering
\begin{tikzpicture}[node distance=0.8cm, scale=0.85, transform shape]
    % Stile file parquet
    \tikzstyle{parquetfile} = [rectangle, rounded corners=2pt, minimum width=2.2cm, minimum height=0.8cm, text centered, draw=black, fill=accentgreen!25, font=\scriptsize]
    
    \node[component, minimum width=3cm] (main) {Main Process};
    
    \node[component, fill=accentgreen!20, below left=1.5cm and 1cm of main] (w1) {Worker 1};
    \node[component, fill=accentgreen!20, below=1.5cm of main] (w2) {Worker 2};
    \node[component, fill=accentgreen!20, below right=1.5cm and 1cm of main] (w3) {Worker N};
    
    \node[above right=0.3cm and -0.5cm of w2] {...};
    
    \draw[arrow] (main) -- (w1);
    \draw[arrow] (main) -- (w2);
    \draw[arrow] (main) -- (w3);
    
    \node[parquetfile, below=1cm of w1] (p1) {file1.parquet};
    \node[parquetfile, below=1cm of w2] (p2) {file2.parquet};
    \node[parquetfile, below=1cm of w3] (p3) {fileN.parquet};
    
    \draw[arrow] (w1) -- (p1);
    \draw[arrow] (w2) -- (p2);
    \draw[arrow] (w3) -- (p3);
\end{tikzpicture}
\caption{Modello fork-join per elaborazione parallela}
\end{figure}

\section{Batching e Scrittura}

Per ottimizzare le scritture su disco, i record vengono accumulati in batch di 10.000 elementi prima del flush:

\begin{lstlisting}[language=Python, caption=Flush batch su Parquet]
BATCH_SIZE = 10000

if len(batch_aiuti) >= BATCH_SIZE:
    pq.write_to_dataset(
        table,
        root_path=str(base_path / "aiuti"),
        partition_cols=['ANNO'],
        existing_data_behavior='overwrite_or_ignore'
    )
\end{lstlisting}

Il parametro \texttt{partition\_cols=['ANNO']} abilita il \textbf{Hive-style partitioning}, creando automaticamente directory \texttt{ANNO=2022/}, \texttt{ANNO=2023/}, etc.
